%=====================================================================
\chapter{Supervised Learning I: Regression}\label{ch:regression}
%=====================================================================
\locator{3}{Part III --- Machine Learning Core}

\begin{outcomes}
By the end of this chapter you will be able to:
\begin{tight}
  \item \textbf{Fit} and \textbf{interpret} a simple and a multiple linear regression, including the meaning of each coefficient.
  \item \textbf{Compute} and \textbf{interpret} MAE, RMSE and $R^2$, and choose between them.
  \item \textbf{Check} the assumptions of linear regression using a residual plot.
  \item \textbf{Explain} how ridge and lasso regularisation control overfitting, and what distinguishes them.
  \item \textbf{Recognise} when a linear model is the wrong tool.
\end{tight}
\end{outcomes}

\begin{prereq}
\chref{math} (gradient descent, dot products), \chref{mlfundamentals}
(train/test, bias--variance) and \chref{eda} (correlation, multicollinearity).
\end{prereq}

\begin{keyterms}
regression $\bullet$ intercept $\bullet$ coefficient $\bullet$ residual $\bullet$
least squares $\bullet$ MAE $\bullet$ MSE $\bullet$ RMSE $\bullet$ $R^2$ $\bullet$
homoscedasticity $\bullet$ polynomial regression $\bullet$ regularisation
$\bullet$ ridge (L2) $\bullet$ lasso (L1)
\end{keyterms}

\section{The Simplest Useful Model}

\begin{definitionbox}[Simple Linear Regression]
Predict a continuous target $y$ from one feature $x$ using a straight line:
\[
\hat{y} = w_0 + w_1 x
\]
where $w_0$ is the \term{intercept} (the prediction when $x=0$) and $w_1$ is the
\term{slope} (the change in $\hat{y}$ per one-unit change in $x$).
\end{definitionbox}

\begin{definitionbox}[Residual and Least Squares]
The \term{residual} for observation $i$ is $e_i = y_i - \hat{y}_i$ --- the vertical
gap between the point and the line. \term{Least squares} chooses $w_0$ and $w_1$ to
minimise $\sum_i e_i^2$.
\end{definitionbox}

\dsfig{%
\begin{tikzpicture}
\begin{axis}[
  width=11.8cm, height=6.0cm,
  xlabel={\scriptsize attendance (\%)}, ylabel={\scriptsize final mark},
  xmin=25, xmax=100, ymin=25, ymax=95,
  axis lines=left, grid=major, grid style={gray!18},
  tick label style={font=\tiny}, label style={font=\scriptsize},
  legend style={font=\tiny, draw=gray!40, at={(0.02,0.98)}, anchor=north west},
]
% fitted line: yhat = 22 + 0.62 x
\addplot[primary, line width=1.4pt, domain=28:98] {22 + 0.62*x};
\addlegendentry{$\hat{y} = 22 + 0.62x$}

% residual segments
\foreach \x/\y in {35/52, 44/44, 52/60, 58/50, 66/70, 72/58, 80/78, 86/68, 92/86, 95/74}{
  \pgfmathsetmacro{\yh}{22+0.62*\x}
  \edef\tmp{\noexpand\draw[accent!70, line width=0.7pt, dashed]
            (axis cs:\x,\y) -- (axis cs:\x,\yh);}\tmp}

\addplot[only marks, mark=*, mark size=1.9pt, primary!80] coordinates {
 (35,52)(44,44)(52,60)(58,50)(66,70)(72,58)(80,78)(86,68)(92,86)(95,74)};
\addlegendentry{observations}

\node[font=\tiny, accent!80!black, anchor=west] at (axis cs:59,44)
     {residual $e_i = y_i - \hat{y}_i$};
\draw[-{Stealth[length=1.6mm]}, accent!70] (axis cs:58.5,45.5) -- (axis cs:58,49);
\node[font=\tiny\itshape, text=gray!55!black, align=center, anchor=north east]
     at (axis cs:99,40) {least squares minimises\\the sum of the \emph{squared}\\dashed lengths};
\end{axis}
\end{tikzpicture}}%
{Simple linear regression. The fitted line minimises the sum of squared residuals ---
squaring, rather than taking absolute values, is what makes large errors
disproportionately costly and gives the closed-form solution.}{regression-line}

\subsection{Interpreting the Coefficients}

\begin{worked}[Reading a Multiple Regression]
A model predicting final mark from three features returns:
\[
\hat{y} \;=\; 18.4 \;+\; 0.41\,(\text{attendance\%}) \;+\; 2.15\,(\text{quizzes}) \;-\; 1.62\,(\text{days\_since\_login})
\]

\smallskip
\textbf{Intercept 18.4:} the predicted mark for a student with 0\% attendance, 0
quizzes and 0 days since last login. This combination is impossible, so the
intercept here is a mathematical anchor, not a meaningful quantity. Report it, do
not interpret it.

\textbf{$+0.41$ on attendance:} each additional percentage point of attendance is
associated with 0.41 more marks, \emph{holding quizzes and recency constant}. Those
five words are the whole meaning of a multiple-regression coefficient.

\textbf{$+2.15$ on quizzes:} each extra quiz completed is worth 2.15 marks, again
holding the others fixed. Note this coefficient is larger than attendance's, but
that does \emph{not} make quizzes more important --- quizzes range 0--10 while
attendance ranges 0--100. Over their full ranges, attendance contributes up to 41
marks and quizzes up to 21.5.

\textbf{$-1.62$ on days since login:} each day of absence costs 1.62 marks. The
negative sign is a sanity check that passed.

\smallskip
\textbf{The causal warning.} None of this says that forcing a student to attend
would raise their mark by 0.41 per point. These are associations, conditional on
the other variables in the model (\chref{causality}).
\end{worked}

\begin{alertbox}[Coefficient Size Is Not Importance]
To compare importance across features, either standardise the features first
(\chref{wrangling}) so all coefficients are on the same scale, or report the effect
over each feature's realistic range. Comparing raw coefficients on different units
is meaningless and is a standard exam trap.
\end{alertbox}

\section{Measuring Regression Error}

\begin{definitionbox}[Three Error Metrics]
\[
\text{MAE} = \frac{1}{n}\sum_{i}\left|y_i - \hat{y}_i\right|
\qquad
\text{RMSE} = \sqrt{\frac{1}{n}\sum_{i}\left(y_i - \hat{y}_i\right)^2}
\qquad
R^2 = 1 - \frac{\sum_i (y_i-\hat{y}_i)^2}{\sum_i (y_i-\bar{y})^2}
\]
MAE and RMSE are in the units of $y$; $R^2$ is the proportion of variance explained
and is unitless.
\end{definitionbox}

\begin{worked}[MAE versus RMSE, and Why the Choice Matters]
Two models predict resolution time (minutes) for five tickets. True values:
$10, 12, 15, 20, 25$.

\smallskip
\textbf{Model A errors:} $3, 3, 3, 3, 3$ \quad (consistently a little off)\\
\textbf{Model B errors:} $0, 0, 0, 0, 15$ \quad (perfect except one disaster)

\smallskip
\textbf{MAE:} A $= 15/5 = 3.0$. \quad B $= 15/5 = 3.0$. \emph{Identical.}

\textbf{RMSE:} A $= \sqrt{45/5} = \sqrt{9} = 3.0$. \quad
B $= \sqrt{225/5} = \sqrt{45} \approx 6.7$.

\smallskip
\textbf{Interpretation.} MAE says the two models are equally good. RMSE says B is
more than twice as bad, because squaring punishes the single 15-minute error far
more than five 3-minute ones.

\textbf{Which is right?} It depends entirely on the cost of a large error:
\begin{tight}
  \item \emph{Ticket triage:} a 15-minute miss on one ticket is no worse than three
        5-minute misses. Use MAE.
  \item \emph{Predicting remaining life of a turbine:} one catastrophic
        underestimate is far worse than many small ones. Use RMSE.
\end{tight}
Choosing the metric is a business decision, made at framing time
(\chref{lifecycle}), not a technical default.
\end{worked}

\section{Checking the Assumptions}

Linear regression makes assumptions, and the residual plot checks nearly all of
them at once. Producing it should be automatic.

\dsfig{%
\begin{tikzpicture}
\begin{groupplot}[
  group style={group size=3 by 1, horizontal sep=0.9cm},
  width=5.35cm, height=3.9cm,
  axis lines=middle, tick label style={font=\tiny}, label style={font=\tiny},
  title style={font=\scriptsize\bfseries},
  xmin=0, xmax=10, ymin=-4.2, ymax=4.2, xtick=\empty, ytick=\empty,
  xlabel={$\hat{y}$}, ylabel={residual},
]
\nextgroupplot[title={\color{greenhl}Good: random scatter}]
\addplot[only marks, mark=*, mark size=1pt, greenhl!80] coordinates {
 (0.8,1.2)(1.5,-0.9)(2.2,0.4)(2.9,-1.6)(3.6,1.8)(4.3,-0.3)(5.0,0.9)(5.7,-1.9)
 (6.4,1.4)(7.1,-0.6)(7.8,0.2)(8.5,-1.3)(9.2,1.0)(1.9,-2.0)(6.8,2.1)(4.7,-2.2)};
\draw[gray!60, dashed] (axis cs:0,0) -- (axis cs:10,0);

\nextgroupplot[title={\color{accent}Bad: curved pattern}]
\addplot[only marks, mark=*, mark size=1pt, accent!80] coordinates {
 (0.8,2.6)(1.5,1.6)(2.2,0.6)(2.9,-0.4)(3.6,-1.4)(4.3,-2.1)(5.0,-2.4)(5.7,-2.2)
 (6.4,-1.5)(7.1,-0.5)(7.8,0.7)(8.5,1.9)(9.2,3.0)};
\draw[gray!60, dashed] (axis cs:0,0) -- (axis cs:10,0);
\node[font=\tiny, accent!80!black, align=center] at (axis cs:5,3.4) {relationship is\\not linear};

\nextgroupplot[title={\color{goldhl!60!black}Bad: fanning out}]
\addplot[only marks, mark=*, mark size=1pt, goldhl] coordinates {
 (0.8,0.2)(1.5,-0.3)(2.2,0.4)(2.9,-0.5)(3.6,0.9)(4.3,-1.0)(5.0,1.4)(5.7,-1.5)
 (6.4,2.1)(7.1,-2.2)(7.8,2.8)(8.5,-3.0)(9.2,3.6)(9.5,-3.5)};
\draw[gray!60, dashed] (axis cs:0,0) -- (axis cs:10,0);
\node[font=\tiny, goldhl!50!black, align=center] at (axis cs:3.2,3.3) {heteroscedastic:\\error grows with $\hat{y}$};
\end{groupplot}
\end{tikzpicture}}%
{Residual plots. Only the left panel is acceptable. A curve means the relationship is
non-linear; a fan means the error variance changes with the prediction, so your confidence
intervals are wrong even if the coefficients are not.}{residual-plots}

\rowcolors{2}{white}{lightbg}
\begin{longtable}{@{}p{3.3cm}p{4.8cm}p{6.2cm}@{}}
\toprule
\rowcolor{primary!15}
\textbf{Assumption} & \textbf{Check} & \textbf{If violated} \\
\midrule
\endhead
Linearity & Residuals vs.\ fitted show no pattern & Add polynomial terms, transform, or use a tree-based model \\
Independence & Domain knowledge; residuals vs.\ time & Use time-series methods (\chref{timeseries}) \\
Homoscedasticity & Residual spread constant across $\hat{y}$ & Log-transform $y$, or use weighted least squares \\
Normal residuals & Q--Q plot or histogram of residuals & Only matters for confidence intervals, not for prediction \\
No multicollinearity & Correlation matrix; variance inflation factor & Drop one of the pair, or use ridge \\
\bottomrule
\end{longtable}

\section{Polynomial Regression and Overfitting}

Adding powers of $x$ lets a linear model fit curves --- $\hat{y} = w_0 + w_1 x +
w_2 x^2 + \dots$ is still \emph{linear in the parameters}, so the same machinery
applies. It is also the fastest way to overfit ever devised.

\section{Regularisation: Ridge and Lasso}

\begin{definitionbox}[Ridge (L2) and Lasso (L1)]
Both add a penalty on coefficient size to the loss:
\[
\text{Ridge: } \sum_i e_i^2 + \lambda \sum_j w_j^2
\qquad\qquad
\text{Lasso: } \sum_i e_i^2 + \lambda \sum_j \left|w_j\right|
\]
$\lambda$ controls the strength. Ridge \emph{shrinks} coefficients towards zero;
lasso can set them \emph{exactly} to zero, performing feature selection.
\end{definitionbox}

\dsfig{%
\begin{tikzpicture}
\begin{groupplot}[
  group style={group size=2 by 1, horizontal sep=1.45cm},
  width=6.6cm, height=4.6cm,
  axis lines=left, tick label style={font=\tiny}, label style={font=\tiny},
  title style={font=\scriptsize\bfseries},
  legend style={font=\tiny, draw=gray!40},
]
\nextgroupplot[title={Ridge: coefficients shrink together},
  xlabel={$\lambda \longrightarrow$}, ylabel={coefficient value},
  xmin=0, xmax=10, ymin=-1.2, ymax=3.2, xtick=\empty,
  legend style={at={(0.99,0.99)}, anchor=north east}]
\addplot[primary, line width=1.1pt, domain=0:10, samples=60]  {2.9/(1+0.42*x)};
\addplot[greenhl, line width=1.1pt, domain=0:10, samples=60]  {1.7/(1+0.42*x)};
\addplot[goldhl, line width=1.1pt, domain=0:10, samples=60]   {0.8/(1+0.42*x)};
\addplot[accent, line width=1.1pt, domain=0:10, samples=60]   {-0.9/(1+0.42*x)};
\draw[gray!55, dashed] (axis cs:0,0) -- (axis cs:10,0);
\node[font=\tiny, text=gray!55!black, align=center, anchor=south east] at (axis cs:9.8,1.6)
     {approach zero,\\never reach it};

\nextgroupplot[title={Lasso: coefficients hit zero and stay},
  xlabel={$\lambda \longrightarrow$},
  xmin=0, xmax=10, ymin=-1.2, ymax=3.2, xtick=\empty]
\addplot[primary, line width=1.1pt] coordinates {(0,2.9)(10,1.05)};
\addplot[greenhl, line width=1.1pt] coordinates {(0,1.7)(6.4,0)(10,0)};
\addplot[goldhl, line width=1.1pt]  coordinates {(0,0.8)(3.1,0)(10,0)};
\addplot[accent, line width=1.1pt]  coordinates {(0,-0.9)(4.6,0)(10,0)};
\draw[gray!55, dashed] (axis cs:0,0) -- (axis cs:10,0);
\addplot[only marks, mark=*, mark size=1.5pt, black] coordinates {(3.1,0)(4.6,0)(6.4,0)};
\node[font=\tiny, text=gray!55!black, align=center, anchor=south east] at (axis cs:9.8,1.7)
     {features are\\removed entirely};
\end{groupplot}
\end{tikzpicture}}%
{Ridge versus lasso as the penalty strength grows. Ridge keeps every feature but makes
each less influential; lasso eliminates features one by one, which is why it doubles as
automatic feature selection.}{ridge-lasso}

\begin{conceptbox}[Which One to Reach For]
\begin{tight}
  \item \textbf{Ridge} when features are correlated and you believe most of them
        matter a little. It handles multicollinearity gracefully.
  \item \textbf{Lasso} when you suspect most features are irrelevant and want a
        short, explainable model. Common with wide clickstream or log-derived
        feature sets.
  \item \textbf{Elastic net} (both penalties) when you cannot decide, which is
        often.
\end{tight}
Regularisation requires scaled features (\chref{wrangling}) --- otherwise the
penalty falls hardest on whichever feature happens to have small units.
\end{conceptbox}

\begin{pitfall}
\begin{tight}
  \item \textbf{Interpreting the intercept} when $x=0$ is impossible.
  \item \textbf{Comparing raw coefficients} across features on different scales.
  \item \textbf{Reporting $R^2$ alone.} $R^2$ never decreases when you add a
        feature, even a random one. Use adjusted $R^2$, or better, validation RMSE.
  \item \textbf{Extrapolating.} A model fitted on attendance 30--95\% says nothing
        about 5\%.
  \item \textbf{Regularising unscaled features.} The penalty becomes arbitrary.
\end{tight}
\end{pitfall}

%---------------------------------------------------------------------
\begin{fourlenses}{Regression}
\lensLA{Predict final mark from week-4 engagement. Report RMSE in marks --- a
figure a programme director understands --- rather than $R^2$. Expect $R^2$ around
0.35--0.5: human outcomes are genuinely noisy, and a claimed 0.95 almost certainly
indicates leakage. Lasso is useful for reducing 300 clickstream features to the
handful you can defend in a meeting.}

\lensPM{Predict task effort in hours from historical features. Use MAE, because
project managers think in ``typically wrong by about 4 hours'' and because effort
data is long-tailed enough that RMSE would be dominated by two disastrous
projects. Beware heteroscedasticity: large tasks have much larger absolute errors,
so consider modelling $\log(\text{hours})$ instead.}

\lensIOT{Predict remaining useful life of a component from sensor trends. Here
RMSE is right, since one large underestimate means an unplanned stoppage. Check
the residual plot for curvature --- degradation is usually non-linear, accelerating
near failure, so a straight line will systematically over-predict life at exactly
the moment accuracy matters most.}

\lensSEC{Regression is less common here than classification, but it appears in
capacity and risk forecasting: predicting expected alert volume for staffing, or
expected breach cost for insurance. Residuals will fan out badly (heteroscedastic)
because incident cost is long-tailed --- model the log, and report a prediction
interval rather than a point estimate.}
\end{fourlenses}

%---------------------------------------------------------------------
\begin{lab}[Fit, Diagnose, Regularise]
\begin{lstlisting}[language=Python]
import numpy as np, matplotlib.pyplot as plt
from sklearn.linear_model import LinearRegression, RidgeCV, LassoCV
from sklearn.metrics import mean_absolute_error, mean_squared_error, r2_score
from sklearn.pipeline import make_pipeline
from sklearn.preprocessing import StandardScaler

# --- fit ---------------------------------------------------------------
lin = LinearRegression().fit(X_tr, y_tr)
pred = lin.predict(X_te)

print("MAE :", mean_absolute_error(y_te, pred))          # typical error
print("RMSE:", mean_squared_error(y_te, pred) ** 0.5)    # penalises big misses
print("R2  :", r2_score(y_te, pred))

for name, coef in zip(feature_names, lin.coef_):
    print(f"  {name:24s} {coef:+.3f}")

# --- diagnose: ALWAYS plot residuals vs fitted -------------------------
plt.scatter(pred, y_te - pred, s=12)
plt.axhline(0, ls="--", c="grey")
plt.xlabel("predicted"); plt.ylabel("residual"); plt.show()
# curve  -> non-linear;  fan -> heteroscedastic;  random cloud -> fine

# --- regularise (note: scaling is mandatory) ---------------------------
ridge = make_pipeline(StandardScaler(), RidgeCV(alphas=np.logspace(-3, 3, 25)))
lasso = make_pipeline(StandardScaler(), LassoCV(cv=5, random_state=0))
ridge.fit(X_tr, y_tr); lasso.fit(X_tr, y_tr)

kept = (lasso[-1].coef_ != 0).sum()
print(f"lasso kept {kept} of {X_tr.shape[1]} features")
\end{lstlisting}
\end{lab}

%---------------------------------------------------------------------
\begin{checkpoint}
\begin{tightnum}
  \item A model has MAE 4.0 and RMSE 11.0. What does the gap between them tell you
        about the error distribution?
  \item In $\hat{y} = 30 + 0.5x_1 + 12x_2$, is $x_2$ twelve times more important
        than $x_1$? Explain.
  \item Your residual plot shows a clear fan shape widening to the right. Name the
        violation and give one remedy.
\end{tightnum}
\end{checkpoint}

%---------------------------------------------------------------------
\begin{chaptersummary}
\begin{tight}
  \item Linear regression fits $\hat{y} = w_0 + \sum_j w_j x_j$ by minimising
        squared residuals.
  \item A coefficient is the change in $\hat{y}$ per unit of that feature
        \emph{holding the others constant} --- and is an association, not a causal
        effect.
  \item MAE treats all errors equally; RMSE punishes large ones. Choose from the
        cost of being badly wrong.
  \item $R^2$ never falls when features are added, so never use it alone for model
        selection.
  \item The residual plot checks linearity and homoscedasticity at a glance.
        Produce it every time.
  \item Ridge shrinks coefficients; lasso zeroes them and thereby selects features.
        Both require scaled inputs.
\end{tight}
\end{chaptersummary}

\begin{reviewq}
  \item \textbf{(MCQ)} Which metric most heavily penalises a single very large
        error? (a)~MAE (b)~RMSE (c)~$R^2$ (d)~median absolute error.
  \item \textbf{(MCQ)} Lasso differs from ridge principally because it can:
        (a)~handle non-linearity (b)~set coefficients exactly to zero
        (c)~work without scaling (d)~fit faster.
  \item \textbf{(Short)} Explain why $R^2$ should not be used to decide whether to
        add a feature.
  \item \textbf{(Short)} A residual plot shows a clear parabola. What is wrong and
        what are two options for fixing it?
  \item \textbf{(Short)} Why must features be standardised before ridge or lasso?
  \item \textbf{(Applied)} Given $\hat{y} = 12.0 + 0.55\,(\text{attendance}) +
        1.9\,(\text{quizzes}) - 0.8\,(\text{days\_absent})$, predict the mark for a
        student with attendance 70\%, 6 quizzes, 3 days absent. Then state, in one
        sentence each, the correct and the incorrect interpretation of the
        attendance coefficient.
  \item \textbf{(Applied)} You are predicting remaining useful life for turbine
        components. Argue for a specific error metric, predict what the residual
        plot will look like and why, and say what you would do about it.
\end{reviewq}
